\color{unidarkgreen}{\subsection{Логика отключения выключателя}}
\color{black}

\begin{enumerate}

\item 
Для вывода воздействия на выключатель при срабатывании функций РЗА предусматривается отдельная логика отключения для каждого выключателя стороны ВН и НН(СН). 

\item
Логика отключения выключателя В2(ОВ)ВН аналогична ЛО В1(ЛВ)ВН, далее приведено описание для ЛО В1(ЛВ)ВН.

\item
Логика отключения выключателя В1(ЛВ) ВН выводится накладкой \textbf{«ЛО В1(ЛВ)ВН»}. 

\item
Сигнал на отключение В1(ЛВ) ВН \textit{«ЛО В ВН: Отключение»} формируется при возникновении любого из следующих событий:
\begin{itemize}
\item прием сигнала срабатывания основных защит;
\item прием сигнала срабатывания резервных защит;
\item прием сигналов срабатывания от внешних защит.
\end{itemize}

\item
Сигнал на запрет АПВ В1(ЛВ) ВН \textit{«ЛО В ВН: Запрет АПВ} формируется при возникновении любого из следующих событий:
\begin{itemize}
\item прием сигнала срабатывания основных защит;
\item прием сигнала срабатывания резервных защит;
\item прием сигналов срабатывания от АПТ и внешних защит стороны НН;
\item прием сигнала запрета АПВ от внешних защит.
\end{itemize}

\item
Сигнал срабатывания от ПО УРОВ НН \textit{«ЛО В ВН: ПО УРОВ} формируется при приеме сигналов срабатывания ПО УРОВ от защит стороны НН.

\item
Алгоритм работы ЛО В1(ЛВ)ВН приведен на рисунке \ref{pic:LO_Q1_LV_VN}. Алгоритм работы ЛО В2(ОВ)ВН выполнен аналогично. Параметры для настройки ЛО В1(ЛВ)ВН приведены в таблице \ref{app_set:lovn1}. Параметры для настройки ЛО В2(ОВ)ВН приведены в таблице \ref{app_set:lovn2}.

\medskip
\begin{figure}[H]
\centering
\includegraphics[width=1\textwidth,height=1\textheight,keepaspectratio]{\fbpath/06.01 ДЗТ/761.9901 ЛО В/img/LO_Q1_LV_VN.pdf}
\captionof{figure}{Функциональная схема логики отключения В1(ЛВ)ВН}\label{pic:LO_Q1_LV_VN}
\end{figure}

\item
Логика отключения выключателей В2 НН1(НН), В1 НН2(СН) и В2 НН2(СН) аналогична ЛО В1 НН1(НН), далее приведено описание для ЛО В1 НН1(НН).

\item
Логика отключения В1 НН1(НН) выводится накладкой \textbf{«ЛО В1 НН1(НН)»}.

\item
Сигнал на отключение В1 НН1(НН) \textit{«ЛО В НН(СН): Отключение»} формируется при возникновении любого из следующих событий:
\begin{itemize}
\item прием сигнала срабатывания основных защит;
\item прием сигнала срабатывания резервных защит;
\item прием сигналов срабатывания от внешних защит;
\item прием сигнала отключения от логики деления;
\item прием сигналов срабатывания от УРОВ и ЗДЗ секций НН;
\end{itemize}

\item
Сигнал на запрет АПВ В1 НН1(НН) \textit{«ЛО В НН(СН): Запрет АПВ»} формируется при возникновении любого из следующих событий:
\begin{itemize}
\item прием сигнала срабатывания основных защит;
\item прием сигнала срабатывания резервных защит;
\item прием сигналов срабатывания от внешних защит;
\item прием сигнала запрета АПВ от внешних защит;
\end{itemize}

\item
Сигнал на запрет АВР секции НН1(НН) \textit{«ЛО В НН(СН): Запрет АВР»} формируется при возникновении любого из следующих событий:
\begin{itemize}
\item прием сигнала срабатывания основных защит;
\item прием сигнала срабатывания резервных защит;
\item прием сигналов срабатывания от внешних защит;
\item прием сигнала отключения от логики деления;
\item прием сигналов срабатывания от УРОВ и ЗДЗ секций НН;
\end{itemize}

\item
Алгоритм работы ЛО В1 НН1(НН) приведен на рисунке \ref{pic:LO_Q1_NN2_SN}. Алгоритм работы ЛО В2 НН1(НН), ЛО В1 НН2(СН), ЛО В2 НН2(СН) выполнен аналогично. Параметры для настройки ЛО В1 НН1(НН) приведены в таблице \ref{app_set:lonn1}. Параметры для настройки ЛО В2 НН1(НН) приведены в таблице \ref{app_set:lonn2}. Параметры для настройки ЛО В1 НН2(СН) приведены в таблице \ref{app_set:lonn3}. Параметры для настройки ЛО В2 НН2(СН) приведены в таблице \ref{app_set:lonn4}.

\medskip
\begin{figure}[H]
\centering
\includegraphics[width=1\textwidth,height=1\textheight,keepaspectratio]{\fbpath/06.01 ДЗТ/761.9901 ЛО В/img/LO_Q1_NN2_SN.pdf}
\captionof{figure}{Функциональная схема логики отключения В1 НН2(СН)}\label{pic:LO_Q1_NN2_SN}
\end{figure}


\end{enumerate}